% ============================================================================
% CAPÍTULO 1 — INTRODUÇÃO
% ============================================================================

\chapter{Introdução}

\section{Contextualização e Tema}

Os genomas mitocondriais --- também denominados \textbf{mitogenomas} --- constituem o conjunto completo de material genético contido na mitocôndria, \textbf{organela} (estrutura subcelular com função especializada) responsável pela produção de energia nas \textbf{células eucarióticas} (células que possuem núcleo delimitado por membrana --- categoria que inclui animais, plantas e fungos, em oposição às bactérias). Essas moléculas de DNA ocupam uma posição central na biologia evolutiva, genética de populações e sistemática, uma vez que oferecem informações valiosas sobre parentesco, variação intra e interespecífica e processos de conservação. O DNA mitocondrial, por suas características estruturais e funcionais --- como a \textbf{herança predominantemente materna} (transmissão exclusiva pela linhagem feminina, sem contribuição paterna), a ausência de \textbf{recombinação} (processo em que \textbf{cromossomos}, as estruturas lineares de DNA no núcleo celular, trocam segmentos durante a reprodução sexual --- ausente no mtDNA, cada molécula é uma cópia integral da mãe) e a elevada taxa evolutiva em comparação ao DNA nuclear ---, tornou-se um marcador molecular amplamente utilizado em diferentes contextos biológicos \cite{boore1999,ekblom2014}.

O avanço das tecnologias de sequenciamento de nova geração (\textit{Next-Generation Sequencing} --- NGS) ampliou o acesso a grandes volumes de dados, favorecendo a obtenção de mitogenomas completos e de alta qualidade. No entanto, a reconstrução dessas moléculas a partir de dados brutos permanece um desafio técnico, especialmente no caso de \textbf{leituras curtas} (\textit{reads} --- os fragmentos individuais de sequência de DNA produzidos pelo sequenciador, tipicamente com 100 a 300 pares de bases), que exigem algoritmos sofisticados para superar problemas como regiões repetitivas e possíveis falhas de \textbf{montagem} --- o processo computacional de reconstruir a sequência original do genoma a partir desses fragmentos sobrepostos. Nesse cenário, ferramentas especializadas, como o NOVOPlasty, têm desempenhado papel fundamental ao implementar estratégias de \textit{seed-and-extend} capazes de recuperar organelas de forma eficiente a partir de dados de genoma total \cite{dierckxsens2017}.

Embora a diversidade de softwares tenha contribuído para avanços significativos na montagem de organelas, a ausência de padronização metodológica ainda compromete a reprodutibilidade das análises. A multiplicidade de versões, dependências e configurações de ambiente gera obstáculos práticos para replicar experimentos, fenômeno amplamente conhecido como \textit{dependency hell} \cite{gruning2018}. Esse problema se insere no contexto mais amplo da crise de reprodutibilidade científica \cite{baker2016}, que tem motivado a adoção, na bioinformática, de práticas da engenharia de software, como a conteinerização de ambientes computacionais com Docker \cite{merkel2014} e a orquestração de fluxos com sistemas de gerenciamento de workflows, a exemplo do Nextflow \cite{ditommaso2017}.

Além da montagem propriamente dita, a anotação funcional dos genes mitocondriais constitui uma etapa igualmente relevante, que permite identificar os genes codificadores de proteínas, RNAs transportadores e RNAs ribossômicos presentes na molécula reconstruída. Ferramentas como o MITOS2 \cite{donath2019} automatizam essa tarefa, mas sua execução --- normalmente restrita a interfaces web --- permanece desvinculada dos pipelines de montagem, criando uma descontinuidade no fluxo analítico.

Soma-se a esses desafios computacionais uma barreira de infraestrutura frequentemente subestimada. Datasets de sequenciamento de genoma total alcançam dezenas a centenas de gigabytes --- o dataset da arara-azul-de-lear (\textit{Anodorhynchus leari}) utilizado neste trabalho, por exemplo, totaliza 87 GB em formato FASTQ descomprimido. Processar esse volume na íntegra exigiria centenas de gigabytes de armazenamento temporário e memória RAM proporcional, recursos indisponíveis em computadores pessoais típicos de estudantes e pequenos laboratórios. Sem estratégias inteligentes de redução de dados, a montagem de mitogenomas a partir de dados públicos ficaria restrita a grupos com acesso a servidores dedicados.

\section{Problema de Pesquisa}

Diante do cenário apresentado --- marcado pela fragmentação das ferramentas de montagem e anotação, pela ausência de padronização metodológica que compromete a reprodutibilidade, pela desconexão entre etapas analíticas que obriga o uso de interfaces web manuais, e pela barreira computacional imposta pelo volume dos dados de sequenciamento ---, este trabalho busca responder à seguinte questão:

\textbf{De que maneira é possível integrar, em um único fluxo de trabalho automatizado, as etapas de aquisição, controle de qualidade, montagem \textit{de novo}, anotação funcional e preparação para submissão ao GenBank de genomas mitocondriais, garantindo reprodutibilidade, portabilidade e viabilidade computacional em hardware doméstico?}

\section{Objetivos}

\subsection{Objetivo Geral}

Desenvolver e validar um pipeline de bioinformática automatizado, reprodutível e portável para a montagem \textit{de novo} e anotação funcional de genomas mitocondriais, utilizando tecnologias de conteinerização (Docker) e sistemas de gerenciamento de workflows (Nextflow).

\subsection{Objetivos Específicos}

\begin{enumerate}[label=\alph*)]
    \item Containerizar, por meio de Dockerfiles com versões pinadas, cada uma das ferramentas bioinformáticas necessárias ao pipeline (SRA-Toolkit, FastQC, Trim Galore, NOVOPlasty e MITOS2).
    \item Desenvolver o script de orquestração do pipeline utilizando Nextflow DSL2, definindo os processos, suas entradas, saídas e interdependências em módulos independentes.
    \item Integrar os containers Docker ao workflow, garantindo que cada etapa da análise seja executada em seu próprio ambiente isolado e pré-configurado.
    \item Implementar um módulo de análise piloto de qualidade (Pilot QC) capaz de determinar automaticamente o volume ideal de dados a serem processados, com base na qualidade das leituras e na fração mitocondrial estimada.
    \item Integrar a anotação funcional automatizada (MITOS2) ao pipeline, incluindo a geração de estruturas secundárias de tRNAs e rRNAs via ViennaRNA.
    \item Automatizar a geração de entregáveis em formato GenBank Flat File, com tratamento adequado de particularidades gênicas como o \textit{frameshift} do ND3 em aves.
    \item Validar o pipeline automatizado por meio de estudos de caso com dados públicos de sequenciamento de pelo menos duas espécies de grupos \textbf{taxonômicos} distintos --- isto é, de posições diferentes na classificação biológica hierárquica (espécie, gênero, família, ordem, etc.) ---, comparando os genomas montados e anotados com resultados de referência.
    \item Disponibilizar o pipeline completo e sua documentação em repositório público no GitHub, assegurando sua acessibilidade e reprodutibilidade pela comunidade científica.
\end{enumerate}

\section{Justificativa}

A análise de mitogenomas representa um componente essencial para a biologia evolutiva, a genética de populações e a sistemática, dado que essas moléculas fornecem informações relevantes sobre relações de parentesco, variação intra e interespecífica e processos de conservação. Sua relevância científica foi consolidada em trabalhos seminais que destacaram o papel estrutural e evolutivo dos mitogenomas como marcadores moleculares \cite{boore1999}, em estudos que enfatizaram sua aplicabilidade em investigações de genética de populações e evolução molecular \cite{ekblom2014} e em contribuições mais recentes que demonstraram seu potencial em análises de biodiversidade e conservação por meio de avanços metodológicos em montagem e anotação de organelas \cite{ulianosilva2023}.

Embora a obtenção de mitogenomas completos tenha sido amplamente favorecida pelo advento do NGS, a execução de pipelines de montagem ainda enfrenta desafios técnicos significativos. A diversidade de ferramentas, versões e dependências envolvidas gera um cenário de difícil replicação, conhecido como \textit{dependency hell}, que compromete a reprodutibilidade e a padronização das análises \cite{sandve2013,baker2016,gruning2018}.

Esse problema se torna especialmente relevante quando aplicado ao estudo de espécies ameaçadas de extinção, nas quais a disponibilidade de amostras biológicas é limitada e a confiabilidade dos resultados genômicos é crucial para fundamentar decisões de conservação. É o caso da \textbf{arara-azul-de-lear} (\textit{Anodorhynchus leari}), ave endêmica da Caatinga nordestina do Brasil, classificada como ``Em Perigo'' pela IUCN Red List. Com população estimada em pouco mais de 1.700 indivíduos na natureza \cite{icmbio2022}, a espécie depende de dados genômicos acessíveis e validáveis para subsidiar programas de manejo, reprodução em cativeiro e monitoramento populacional. Embora seu mitogenoma já tenha sido previamente reconstruído em atividade acadêmica, essa montagem foi realizada de forma manual, sem padronização metodológica e sem disponibilização pública do fluxo de trabalho --- o que impede sua verificação e replicação por outros pesquisadores.

Um desafio adicional reside na completude dos pipelines existentes. A maioria das ferramentas e workflows publicados concentra-se exclusivamente na etapa de montagem, sem integrar a anotação funcional nem a preparação dos resultados para submissão a bancos de dados públicos como o GenBank. Essa lacuna obriga o pesquisador a recorrer a interfaces web externas (e.g., MITOS2 via UseGalaxy) e a conversões manuais de formato, introduzindo etapas não rastreáveis que comprometem a reprodutibilidade do ciclo completo de análise.

Nesse contexto, torna-se necessário o desenvolvimento de soluções que conciliem robustez biológica, consistência computacional e acessibilidade de hardware. A adoção de práticas modernas da engenharia de software, como a containerização e os sistemas de gerenciamento de workflows, tem se mostrado fundamental para superar limitações de reprodutibilidade, além de permitir que fluxos analíticos sejam portáveis, escaláveis e transparentes \cite{ditommaso2017,ewels2020}.

Assim, este trabalho justifica-se pela proposta de construir um pipeline automatizado e reprodutível que cubra o ciclo completo --- desde a aquisição dos dados brutos até a geração de arquivos prontos para submissão ao GenBank (banco de dados público mantido pelo NCBI para depósito de sequências genéticas) ---, unindo boas práticas de ciência aberta e engenharia computacional à relevância biológica dos mitogenomas. A aplicação ao caso concreto da arara-azul-de-lear demonstra não apenas a viabilidade técnica da abordagem, mas também sua contribuição direta para a conservação de uma espécie ameaçada da biodiversidade brasileira.

\section{Metodologia}

O presente trabalho adota uma abordagem de pesquisa aplicada, de natureza experimental, com desenvolvimento iterativo e incremental de software. O pipeline foi implementado em Nextflow DSL2 e containerizado com Docker, seguindo boas práticas de engenharia de software científico \cite{wilson2014}.

A validação foi realizada por meio de dois estudos de caso com dados públicos do NCBI Sequence Read Archive (SRA): (i) \textit{Diploprion bifasciatus} (Perciformes --- peixe-sabão barrado), como espécie-controle para verificação funcional do pipeline; e (ii) arara-azul-de-lear (\textit{Anodorhynchus leari}, Psittaciformes --- aves), como espécie-alvo para a aplicação final. Em ambos os casos, os mitogenomas foram montados \textit{de novo} a partir de leituras Illumina de genoma total e comparados com referências disponíveis no GenBank.

O pipeline integra sete etapas sequenciais automatizadas: (1) análise piloto de qualidade para dimensionamento do volume de dados; (2) download e truncamento dos dados do SRA; (3) controle de qualidade com FastQC; (4) remoção de adaptadores e bases de baixa qualidade com Trim Galore; (5) montagem \textit{de novo} com NOVOPlasty; (6) anotação funcional com MITOS2; e (7) compilação automática de 14 categorias de entregáveis, incluindo arquivos em formato GenBank Flat File prontos para submissão. Os detalhes de cada etapa, incluindo ferramentas, parâmetros e versões utilizadas, são apresentados no Capítulo~\ref{cap:metodologia}.

\section{Estrutura do Trabalho}

O presente trabalho está organizado em sete capítulos. Este primeiro capítulo apresentou a contextualização do tema, o problema de pesquisa, os objetivos, a justificativa e um breve resumo metodológico.

O \textbf{Capítulo~\ref{cap:referencial}} reúne o referencial teórico, abordando os fundamentos da montagem \textit{de novo} de genomas, as especificidades do DNA mitocondrial, a crise de reprodutibilidade em ciência computacional, a tecnologia de conteinerização com Docker e os sistemas de gerenciamento de workflows científicos.

O \textbf{Capítulo~\ref{cap:trabalhos}} analisa os trabalhos relacionados, revisando as principais ferramentas disponíveis para montagem de organelas --- MitoHiFi, MITObim, GetOrganelle, MToolBox e NOVOPlasty --- e apresenta uma análise comparativa entre elas.

O \textbf{Capítulo~\ref{cap:metodologia}} detalha a metodologia, descrevendo o processo de desenvolvimento de software adotado, a visão geral e o fluxo de execução do pipeline, sua arquitetura modular, as ferramentas e tecnologias utilizadas, e cada uma das sete etapas da análise.

O \textbf{Capítulo~\ref{cap:resultados}} apresenta os resultados obtidos nas execuções do pipeline e a discussão, incluindo as métricas de montagem e anotação para ambas as espécies, a análise de desempenho computacional e uma avaliação da robustez da abordagem.

O \textbf{Capítulo~\ref{cap:mitogenoma}} é dedicado integralmente à arara-azul-de-lear, apresentando a caracterização detalhada de seu mitogenoma, a análise dos genes codificadores de proteínas, RNAs transportadores e ribossomais, a comparação com a espécie congênere \textit{Anodorhynchus hyacinthinus} e as implicações dos resultados para a conservação da espécie.

O \textbf{Capítulo~\ref{cap:consideracoes}} encerra o trabalho com as considerações finais, sintetizando as contribuições alcançadas e indicando perspectivas para trabalhos futuros.
